% BHU PYQ -- Transcribed & Corrected
% Paper: PHYMJ21/PHYMN21: Thermal Physics
% Exam: B.Sc. (Hons) Semester II (NEP) Examination, 2024-25
% Ref Code: CECONF/N/2024-25/4
% Category: Major / Minor
% Department: Physics

\documentclass[11pt,a4paper]{article}
\usepackage[a4paper,top=2.5cm,bottom=2.5cm,left=2.8cm,right=2.8cm,headheight=14pt]{geometry}
\usepackage{amsmath,amssymb}
\usepackage[T1]{fontenc}
\usepackage[utf8]{inputenc}
\usepackage{lmodern,microtype,fancyhdr,enumitem,hyperref,lastpage,parskip}

\hypersetup{
    colorlinks=true,
    urlcolor=black,
    linkcolor=black,
    pdftitle={PHYMJ21/PHYMN21: Thermal Physics -- BHU Sem II 2024-25 (NEP)}
}

\pagestyle{fancy}
\fancyhf{}
\renewcommand{\headrulewidth}{0.4pt}
\renewcommand{\footrulewidth}{0.4pt}
\fancyhead[L]{\small \textit{Banaras Hindu University}}
\fancyhead[R]{\small \textit{PHYMJ21/PHYMN21: Thermal Physics (2024-25)}}
\fancyfoot[C]{\small Page \thepage\ of \pageref{LastPage}}

\newcommand{\pts}[1]{\hfill \textit{[#1]}}

\begin{document}

\begin{center}
  {\large\bfseries BANARAS HINDU UNIVERSITY}\\[2pt]
  {\normalsize B.Sc. (Hons) Semester II Examination, 2024-25 (NEP)}\\[6pt]
  \rule{0.8\linewidth}{0.5pt}\\[6pt]
  {\large\bfseries PHYSICS}\\[4pt]
  {\large\bfseries Paper Code: PHYMJ21 / PHYMN21 : Thermal Physics}\\[2pt]
  {\small Ref. No.: CECONF/N/2024-25/4}\\[4pt]
  {\normalsize Time: 2\frac{1}{2} Hours\hfill Maximum Marks: 70}
\end{center}

\rule{\linewidth}{0.4pt}

\smallskip

\noindent \textit{Instructions: Attempt five questions including Question 1 which is compulsory. The figures in the right-hand margin indicate marks.}

\bigskip

\noindent \textbf{Question 1.} Answer all questions briefly: \pts{20}
\begin{enumerate}[label=(\alph*)]
  \item State the zeroth law of thermodynamics.
  \item Explain Zartman and Ko's experiment for verification of Maxwell's law of distribution of speeds.
  \item Calculate the mean free path for molecules of hydrogen gas. The radius of the gas molecule is $1.37 \times 10^{-10}\text{ m}$ and the density of gas is $3 \times 10^{25}\text{ m}^{-3}$.
  \item Write photon gas pressure inside a cavity in terms of photon's energy density.
  \item Draw T-S diagram of a reversible Carnot cycle.
  \item Calculate the change in entropy when $1.119\text{ kg}$ of water at $100^\circ\text{C}$ is converted into steam at the same temperature. Latent heat of steam is $540\text{ cal/gram}$.
  \item Write inversion temperature in terms of real gas constants.
  \item Write the relation between Gibbs and Helmholtz free energy.
  \item Explain the triple point of water using a suitable diagram.
  \item Explain Joule-Thomson effect.
\end{enumerate}

\bigskip

\noindent \textbf{Question 2.}
\begin{enumerate}[label=(\alph*)]
  \item Explain the Carnot cycle. Show that no irreversible heat engine operating between two reservoirs (source and sink having absolute temperatures $T_1$ and $T_2$ respectively) can be more efficient than a reversible engine operating between the same reservoirs. \pts{8$\frac{1}{2}$}
  \item Calculate the work done in kJ during an isothermal reversible process when 5 moles of ideal gas are expanded so that its volume is doubled at a temperature of 400 K. \pts{4}
\end{enumerate}

\bigskip

\noindent \textbf{Question 3.}
\begin{enumerate}[label=(\alph*)]
  \item Show that the change in entropy is zero for a reversible cyclic process, but it increases for an irreversible process. \pts{7}
  \item Write the relation between heat capacities at constant pressure and at constant volume as a function of isothermal compressibility and volume expansivity. Calculate the difference between them for $^4\text{He}$ gas at temperature 6 K and pressure 19.7 atm. [\textit{Given:} specific volume $= 2.64 \times 10^{-2}\text{ m}^3\text{ kmol}^{-1}$; isothermal compressibility $= 9.42 \times 10^{-8}\text{ m}^2\text{N}^{-1}$; volume expansivity $= 5.35 \times 10^{-2}\text{ K}^{-1}$]. \pts{5$\frac{1}{2}$}
\end{enumerate}

\bigskip

\noindent \textbf{Question 4.}
\begin{enumerate}[label=(\alph*)]
  \item Derive Clausius-Clapeyron equation using Carnot's cycle. \pts{7$\frac{1}{2}$}
  \item Prove the following relation: \pts{5}
  \[ T \cdot dS = C_p dT - T \left(\frac{\partial V}{\partial T}\right)_p dP \]
\end{enumerate}

\bigskip

\noindent \textbf{Question 5.}
\begin{enumerate}[label=(\alph*)]
  \item Derive the Rayleigh-Jeans law. Also explain the UV catastrophe. \pts{8$\frac{1}{2}$}
  \item The sun's radius is $6.96 \times 10^8\text{ m}$ and its surface temperature is $5700\text{ K}$. Estimate the total power radiated by the Sun. (Stefan's constant $\sigma = 5.67 \times 10^{-8}\text{ W m}^{-2}\text{ K}^{-4}$). \pts{4}
\end{enumerate}

\bigskip

\noindent \textbf{Question 6.}
\begin{enumerate}[label=(\alph*)]
  \item Explain transport phenomena in gases. Derive the expression for viscosity of gases. \pts{7$\frac{1}{2}$}
  \item Write critical gas constants in terms of real gas constants. Find the gas constant $b$ if critical constant $V_c$ is $69.7\text{ cc/mole}$. \pts{5}
\end{enumerate}

\bigskip

\noindent \textbf{Question 7.}
\begin{enumerate}[label=(\alph*)]
  \item Derive the following Maxwell's thermodynamic relation: \pts{5}
  \[ \left(\frac{\partial S}{\partial V}\right)_T = \left(\frac{\partial P}{\partial T}\right)_V \]
  \item Explain the principle of regenerative cooling with a schematic diagram. \pts{7$\frac{1}{2}$}
\end{enumerate}

\bigskip
\begin{center}
  \textbf{***}
\end{center}

\end{document}
